\vspace{-1pt}
\section{Related Work}
The related work will be discussed from two different perspectives. First, we will discuss network sampling bias, and the consequences thereof. Next, we will introduce state-of-the-art community detection approaches with a special focus on those that we use as baselines in this work.

There is a body of work surrounding different types of biases in social networks, such as the retweet network~\cite{Diaz:2018:AAB:3173574.3173986,Liao:2016:SUB:2858036.2858422}. However, these studies do not specify the root of the existing biases in social networks and ways we can mitigate them through community detection. In this paper, we show that the root of existing biases in social networks come from not only the network structure but also community detection methods that exemplify bias by putting lowly-connected nodes into non-significant communities which leads to their ignorance in the downstream studies of these communities. Therefore, we analyze some state of the art community detection methods and investigate their flaws and weaknesses. 

Current community detection methods, such as relying on the modularity value in network structures~\cite{blondel2008fast}, suffer from bias of ignoring lowly-connected users. 
This is a type in which users with non-genuine retweets get assigned to significant communities for significant analysis, while users with low retweet rates get assigned to non-significant groups which leads to their exclusion from the study. On the other hand, some other community detection methods~\cite{6729613} suffer from low recall scores due to their exclusion of users from their belonging communities. Community detection methods with low recall scores can suffer from incorporating bias by excluding so many significant users. Although there are many different methods for the community detection task~\cite{Papadopoulos2012,gargi2011large}, we divided our analysis in two different sections: 
\begin{quote}
\begin{enumerate}
\item Methods that do not use node attributes for the community detection task.
\item Methods that use node attributes for the community detection task.
\end{enumerate}
\end{quote}
Despite the fact that there are different methods that can go under each of these two major groups discussed above, we picked the two most well known and broadly used approaches to conduct our analysis later in this paper.

\subsection{Non-Attributed Methods}
One of the widely used community detection methods is the Louvain method, which utilizes the modularity value to obtain partitions by optimizing an objective function discussed in~\cite{blondel2008fast}. However, this method tends to create small communities by putting introverted and low-degree users into small communities which can be biased towards low-degree people who are relevant by context to the significant groups but not connected through the network structure. The results of such instances will be discussed in detail in the results section of our proposed method, but we will give the reader a brief overview of such cases in this section. 
Placement of low-degree people into insignificant groups of their own may create the illusion that these small communities are communities that are not so relevant to the major groups that need to be studied and therefore can be ignored, but we observed that these communities had so much relevant and informative content to the study, based on our ground truth labels, that it is crucial to keep them in the significant communities. Merely low-degree rates and network structure should not cause a user to be ignored.
Low-connection in social networks may be as a result of various factors. Consider the retweet network as an example, low-connection may have different reasons, such as users having unique ways of expressing ideas yet relevant to the majority group's opinion, the tweet content being long which tends to repel other users to read and retweet them or in case of introverted users with small friendship circles tweets may not get enough attention or retweets. Network sampling can also change network measures and create unrealistic introverted users~\cite{gonzalez2014assessing,morstatter2013sample}. These reasons are not indicative of irrelevancy of topic and should not cause bias towards low-degree users of such types and their exclusion from significant communities.
%Low retweet rates can put these users into small introverted communities due to various reasons, such as the tweet content being long which tends to repel other users to read and retweet them or in case of introverted users with small friendship circles their tweets may not get enough attention or retweet rates, but this does not indicate that the user should not belong to one of the major communities. 
Instead what the modularity value does is to separate these users into introverted and small communities depriving them from being into their correct significant communities by only focusing on the network structure and its connections. Other methods of such type that do not utilize node attributes exist, such as BigCLAM and DEMON~\cite{yang2013overlapping,coscia2012demon}. Using node attributes in addition to the network structure will add additional information, as discussed in \cite{cho2016latent}, and can be helpful to reduce this type of biases in social networks. There are methods that use node attributes in combination with the network structure to create communities. In the next section, we will discuss some methods that utilize node attributes in addition to the network structure and discuss issues associated with them and the potential biases they create.

\subsection{Attributed Methods}
There are some attributed community detection methods discussed in \cite{Falih:2018:CDA:3184558.3191570}. From which, one of the widely known methods that uses node attributes in addition to the network structure is a generative model called CESNA~\cite{6729613}. This method restructures the network by incorporating node attributes. Experiments show that this method performs well by having a high precision value; however, it suffers from having a significantly low recall value. Low recall value may be indicative that this method also discards some of the important users which we are trying to avoid in order to minimize bias. In addition to CESNA, other methods have been proposed, such as Block-LDA~\cite{balasubramanyan2011block}, PAICAN~\cite{Bojchevski2018BayesianRA}, shared latent space models such as CLSM~\cite{cho2016latent}, and embedding based approaches like  LANE~\cite{huang2017label}, which utilize node attributes in order to detect communities, or ELAINE~\cite{Goyal:2018:ENE:3209542.3209571} which utilizes edge attributes. However, none of these methods tried to target the existing bias in social networks.
